
\documentclass[11pt,a4paper]{article}
\usepackage[margin=2.35cm]{geometry}
\usepackage[T1]{fontenc}
\usepackage{lmodern,microtype}
\usepackage{amsmath,amssymb,booktabs,longtable,array}
\usepackage{enumitem,xcolor,hyperref}
\usepackage{algorithm,algpseudocode}
\setlength{\emergencystretch}{2em}
\hypersetup{colorlinks=true,linkcolor=blue!55!black,urlcolor=blue!55!black}
\newcommand{\qt}{\textsc{QUETRANSPORT}}
\newcommand{\Rland}{\mathcal{R}}
\newcolumntype{P}[1]{>{\raggedright\arraybackslash}p{#1}}
\title{\textbf{Codebook for QUETRANSPORT}\\
\large A general quantitative spatial toolkit for transport appraisal}
\author{Gabriel M. Ahlfeldt}
\date{\today}

\begin{document}
\maketitle
\begin{abstract}
\qt\ implements a deliberately streamlined version of the quantitative urban
model in Ahlfeldt, Redding, Sturm, and Wolf (2015, ARSW) for an arbitrary
number of spatial units. It takes standardized spatial data and baseline and
policy travel-time matrices, inverts local fundamentals conditional on chosen
parameters, solves closed- and open-city counterfactual equilibria, and
evaluates a fixed-distribution accounting benchmark. Each is also evaluated
after switching off productivity and amenity spillovers. This
codebook states the assumptions and equations implemented in MATLAB, shows
that the equilibrium has as many equations as unknown target variables,
documents every inverted primitive, and distinguishes the common observed
floor-space rent, endogenous use-specific bid rents, and annual land rent.
Structural estimation is outside the toolkit.
\end{abstract}
\tableofcontents
\newpage

\section{Purpose, scope, and sequence}
The toolkit answers a conditional question: given a quantitative spatial model,
a baseline city, parameter values, and a transport change, what new spatial
equilibrium does the model imply? It does not estimate the structural
parameters. The user chooses variable mappings, parameters, shocks, numerical
controls, and city closure in the root file \texttt{project\_config.yaml}.

The workflow has four economic steps:
\begin{enumerate}
\item specify parameter values and assemble mutually consistent spatial inputs;
\item invert the baseline fundamentals that rationalize the observed allocation;
\item change travel times and, optionally, productivity, amenity, or development
fundamentals;
\item solve the baseline and policy equilibria under the same closure, evaluate
the fixed-distribution benchmark, and report local, aggregate, welfare, and
land-rent effects.
\end{enumerate}
Every location vector has length $N$. Both travel-time matrices are
$N\times N$ and use the same location ordering as the standardized table and
shapefile. GRID may construct regular squares or hexagons, or use
\texttt{cell\_geometry: original}. In the latter case each supplied source
polygon is retained as a model location, the cell-size parameter is ignored,
and the same polygons determine routing centroids, MATLAB locations, shock
interpolation targets, and reporting maps. A configured unique source ID may
be retained as the canonical location identifier.

\section{Economic environment and assumptions}
Locations are indexed by residence $i$ and workplace $j$. The model has the
following maintained assumptions.
\begin{enumerate}[label=\textbf{A\arabic*.}]
\item There is a continuum of workers. Each supplies one unit of labour and
chooses a residence--workplace pair.
\item Workers have Frechet-distributed idiosyncratic preferences. This produces
a smooth gravity equation for commuting and a closed-form expected-utility
index.
\item Commuting reduces the effective wage exponentially in bilateral travel
time. Travel times are measured in the unit used to calibrate $\kappa$.
\item A homogeneous, freely traded final good is the numeraire. Firms are
perfectly competitive and use labour and commercial floor space under constant
returns to scale.
\item Residents spend share $\beta$ of income on the numeraire and share
$1-\beta$ on residential floor space.
\item Geographic land is fixed. A reduced-form developer technology converts
land and the non-land development input into floor space.
\item Productivity and amenities may depend on surrounding employment and
population densities. Setting the corresponding spillover elasticity to zero
turns that feedback off.
\item In a closed city total population is fixed and utility adjusts. In an
open city utility is fixed at the outside option and total population adjusts.
\item The baseline inversion is conditional on parameter values and observed
quantities and rents. It is an exact model-based quantification, not an
estimation exercise.
\end{enumerate}

\section{Workers, commuting, and welfare}
Let $B_i$ be total amenity, $Q_i$ residential floor-space rent, $w_j$ the wage,
and $\tau_{ij}$ travel time. Define bilateral attractiveness
\begin{equation}
 \Phi_{ij}=
 e^{-\epsilon\kappa\tau_{ij}}
 B_i^\epsilon Q_i^{-(1-\beta)\epsilon}w_j^\epsilon,
 \qquad
 \Phi=\sum_i\sum_j\Phi_{ij}.                         \label{eq:phi}
\end{equation}
The unconditional probability of residing in $i$ and working in $j$ is
\begin{equation}
 \pi_{ij}=\frac{\Phi_{ij}}{\Phi}.                     \label{eq:pi}
\end{equation}
For total city population $H$,
\begin{equation}
 H_{Ri}=H\sum_j\pi_{ij},\qquad
 H_{Mj}=H\sum_i\pi_{ij},\qquad
 \sum_iH_{Ri}=\sum_jH_{Mj}=H.                        \label{eq:margins}
\end{equation}
Conditional commuting probabilities and resident income are
\begin{equation}
 \pi_{j|i}=\frac{\Phi_{ij}}{\sum_s\Phi_{is}},\qquad
 \bar w_i=\sum_j\pi_{j|i}w_j,\qquad
 v_i=H_{Ri}\bar w_i.                                  \label{eq:income}
\end{equation}
Thus $v_i$ is total labour income of residents of $i$, which is the object used
in residential floor-space demand. Expected utility is
\begin{equation}
 U=\gamma\Phi^{1/\epsilon},\qquad
 \gamma=\Gamma\!\left(\frac{\epsilon-1}{\epsilon}\right). \label{eq:utility}
\end{equation}

\section{Production, floor space, and land}
Total floor space in location $i$ is $L_i$. Share $\theta_i$ is commercial, so
$L_{Mi}=\theta_iL_i$ and $L_{Ri}=(1-\theta_i)L_i$. Production is
\begin{equation}
 Y_i=A_iH_{Mi}^{\alpha}L_{Mi}^{\,1-\alpha}.           \label{eq:production}
\end{equation}
Perfect competition gives the labour and commercial-floor-space conditions
\begin{equation}
 w_i=\frac{\alpha Y_i}{H_{Mi}},\qquad
 q_i=\frac{(1-\alpha)Y_i}{L_{Mi}},                    \label{eq:firmfoc}
\end{equation}
where $q_i$ is commercial floor-space rent. Residential market clearing is
\begin{equation}
 Q_iL_{Ri}=(1-\beta)v_i,                              \label{eq:resclear}
\end{equation}
where $Q_i$ is residential floor-space rent.

Let $K_i$ denote geographic land and let $\mu$ be the non-land share in
development. The reduced-form floor-space stock is
\begin{equation}
 L_i=\varphi_iK_i^{\,1-\mu}.                          \label{eq:floorsupply}
\end{equation}
The configured values are $\mu=0.75$ and $1-\mu=0.25$. Hence land enters
\eqref{eq:floorsupply} with exponent $0.25$, not $0.75$.

\subsection{Mixed and specialized locations}
The active support of productivity and amenity determines whether a location
can host firms, residents, or both. For a mixed location, no arbitrage between
uses implies
\begin{equation}
 q_i=Q_i=r_i^F,                                       \label{eq:noarb}
\end{equation}
and combining \eqref{eq:firmfoc}--\eqref{eq:resclear} yields
\begin{align}
 r_i^F&=\frac{(1-\alpha)Y_i+(1-\beta)v_i}{L_i},       \label{eq:mixedrent}\\
 \theta_i&=\frac{(1-\alpha)Y_i}{r_i^FL_i}.            \label{eq:mixedshare}
\end{align}
In a commercially specialized location $\theta_i=1$ and the commercial
condition determines $q_i$. In a residentially specialized location
$\theta_i=0$ and the residential condition determines $Q_i$.

\subsection{One observed rent and two endogenous bid rents}
The standardized input supplies one common baseline floor-space rent,
$R_i^{obs}$. The inversion uses this same observation in both the commercial
and residential demand equations. It is neither land rent nor a regulatory
wedge. This is a transparent maintained restriction on baseline data.

The counterfactual model nevertheless retains two endogenous bid rents:
commercial $q_i$ and residential $Q_i$. Economic demand differs across the two
uses. In a mixed-use location the no-arbitrage condition \eqref{eq:noarb}
imposes equality; at a specialized corner only the active use's bid rent clears
its market, so $q_i$ and $Q_i$ need not coincide. Annual land rent
$\Rland_i$ remains the separate developer residual in \eqref{eq:landrent}.

\subsection{Floor-space rents, annual land rent, and land value}
The model contains three distinct concepts:
\begin{enumerate}
\item $q_i$: price per unit of commercial floor space;
\item $Q_i$: price per unit of residential floor space;
\item $\Rland_i$: annual rent paid to the fixed factor land.
\end{enumerate}
Under the constant-returns developer problem, Euler's theorem assigns the land
share $1-\mu$ of gross floor-space revenue to land:
\begin{equation}
 \boxed{\Rland_i=(1-\mu)
 \left(q_iL_{Mi}+Q_iL_{Ri}\right).}                  \label{eq:landrent}
\end{equation}
This is the object computed by
\texttt{qt\_write\_closure\_results.m} and reported as
\texttt{TotalLandRentPct}. It is a \emph{derived endogenous equilibrium
outcome}: rents and land-use shares change in equilibrium, and
\eqref{eq:landrent} maps those outcomes into the developer residual. It is not
a floor-space rent and is not their unweighted sum. A capitalized land asset
value would additionally require a capitalization rate (and expectations about
future rents), so the toolkit does not label annual land rent as land value.

Aggregate annual land rent is
\begin{equation}
 \Rland=\sum_i\Rland_i,\qquad
 \%\Delta\Rland=100\left(\frac{\Rland^1}{\Rland^0}-1\right). \label{eq:aggland}
\end{equation}

\section{Endogenous productivity and amenity}
Employment-density and residential-density spillovers are
\begin{align}
 \Upsilon_i&=\sum_s e^{-\delta\tau_{is}}\frac{H_{Ms}}{K_s},
 & A_i&=a_i\Upsilon_i^\lambda,\\
 \Omega_i&=\sum_r e^{-\rho\tau_{ir}}\frac{H_{Rr}}{K_r},
 & B_i&=b_i\Omega_i^\eta.                             \label{eq:spillovers}
\end{align}
Here $a_i$ and $b_i$ are fundamentals held fixed across a pure transport
counterfactual; $A_i$ and $B_i$ are total productivity and amenity, which
respond endogenously. Setting $\lambda=0$ or $\eta=0$ makes the corresponding
total object equal to its fundamental.

\section{Objects and their status}
\small
\begin{longtable}{P{2.8cm}P{10.7cm}}
\caption{Structural, supplied, inverted, and endogenous objects}\label{tab:objects}\\
\toprule Object & Meaning and status\\ \midrule
\endfirsthead
\multicolumn{2}{l}{\emph{Table \ref{tab:objects} continued}}\\
\toprule Object & Meaning and status\\ \midrule
\endhead
\multicolumn{2}{l}{\textbf{Parameters fixed by the user}}\\
$\alpha,\beta$ & Labour share in production and numeraire-consumption share.\\
$\epsilon,\kappa$ & Commuting-choice elasticity and time-cost coefficient.\\
$\lambda,\delta$ & Productivity-spillover elasticity and spatial decay.\\
$\eta,\rho$ & Amenity-spillover elasticity and spatial decay.\\
$\mu,1-\mu$ & Non-land and land shares in development.\\
\addlinespace
\multicolumn{2}{l}{\textbf{Baseline data or externally supplied objects}}\\
$K_i$ & Geographic land area.\\
$H_{Ri}^{obs}$ & Resident population.\\
$H_{Mi}^{obs}$ & Workplace employment, rescaled to equal population in aggregate.\\
$R_i^{obs}$ & One common observed baseline floor-space rent, used in both inversion conditions.\\
$\tau^0,\tau^1$ & Baseline and counterfactual bilateral travel times.\\
\addlinespace
\multicolumn{2}{l}{\textbf{Primitives recovered by baseline inversion}}\\
$a_i,b_i$ & Productivity and amenity fundamentals net of endogenous spillovers.\\
$\varphi_i$ & Structural-density/floor-supply fundamental.\\
$\bar U$ & Baseline utility used as outside utility in the open-city model.\\
\addlinespace
\multicolumn{2}{l}{\textbf{Adjusted baseline objects recovered during inversion}}\\
$A_i,B_i$ & Total productivity and amenity that rationalize the two employment margins.\\
$w_i,\bar w_i,v_i$ & Workplace wage, expected wage, and resident total income.\\
$L_i,L_{Mi},L_{Ri},\theta_i$ & Total and use-specific floor space and commercial share.\\
\addlinespace
\multicolumn{2}{l}{\textbf{Endogenous counterfactual equilibrium outcomes}}\\
$\pi_{ij},H_{Ri},H_{Mi}$ & Commuting probabilities and residential/workplace employment.\\
$A_i,B_i,Y_i$ & Total productivity, total amenity, and output.\\
$w_i,q_i,Q_i,\theta_i$ & Wage, two endogenous use-specific bid rents, and commercial floor-space share.\\
$U,H$ & Utility and population; which is fixed depends on closure.\\
$\Rland_i,\Rland$ & Annual local and aggregate land rent from the developer residual.\\
\bottomrule
\end{longtable}
\normalsize

\section{Baseline inversion}
The inversion conditions on
$\{R^{obs},H_M^{obs},H_R^{obs},K,\tau^0\}$ and the parameter vector.
Employment is first rescaled so
$\sum_iH_{Mi}^{obs}=\sum_iH_{Ri}^{obs}=H$.

\subsection{Adjusted productivity, amenity, and wages}
Combining the firm's two first-order conditions eliminates commercial floor
space and gives the wage equation implemented in \texttt{cmodexog.m}:
\begin{equation}
 w_i=\alpha A_i^{1/\alpha}
 \left(\frac{1-\alpha}{R_i^{obs}}\right)^{(1-\alpha)/\alpha}. \label{eq:invwage}
\end{equation}
Given guesses for $A$ and $B$, equations
\eqref{eq:phi}--\eqref{eq:margins} produce predicted employment margins
$\widehat H_M$ and $\widehat H_R$. The multiplicative inversion mappings are
\begin{align}
 A_i^{*}&=A_i
 \left(\frac{H_{Mi}^{obs}}{\widehat H_{Mi}}\right)^{1/\epsilon},&
 B_i^{*}&=B_i
 \left(\frac{H_{Ri}^{obs}}{\widehat H_{Ri}}\right)^{1/\epsilon}. \label{eq:abupdate}
\end{align}
The code damps these mappings, normalizes the geometric mean of positive $A_i$
to one, and rescales $B$ so $\Phi=H$. At convergence,
\begin{equation}
 \widehat H_{Mi}=H_{Mi}^{obs},\qquad
 \widehat H_{Ri}=H_{Ri}^{obs}\quad\forall i.            \label{eq:invmargins}
\end{equation}
The productivity normalization fixes the nominal scale; amenity scaling adopts
the ARSW utility normalization.

\subsection{Floor space and structural density}
After convergence, all remaining floor-space objects are recovered directly:
\begin{align}
 L_{Mi}&=\left(\frac{(1-\alpha)A_i}{R_i^{obs}}\right)^{1/\alpha}H_{Mi}^{obs},
                                                               \label{eq:invLM}\\
 L_{Ri}&=\frac{(1-\beta)v_i}{R_i^{obs}},                        \label{eq:invLR}\\
 L_i&=L_{Mi}+L_{Ri},\qquad
 \theta_i=\frac{L_{Mi}}{L_i},                                  \label{eq:invtheta}\\
 \varphi_i&=\frac{L_i}{K_i^{\,1-\mu}}.                          \label{eq:invvarphi}
\end{align}
The same $R_i^{obs}$ enters both use-specific equations, but it must not be
confused with land rent. This is the toolkit's maintained one-rent baseline
restriction; it does not impose equal commercial and residential demand.

\subsection{Productivity and amenity fundamentals}
Using observed employment densities and baseline travel times, compute
$\Upsilon_i^0$ and $\Omega_i^0$ from \eqref{eq:spillovers}. The fundamentals
held fixed in a pure transport experiment are
\begin{equation}
 a_i=\frac{A_i}{(\Upsilon_i^0)^\lambda},\qquad
 b_i=\frac{B_i}{(\Omega_i^0)^\eta}.                        \label{eq:invfund}
\end{equation}
Finally, baseline outside utility follows from \eqref{eq:utility}:
\begin{equation}
 \bar U=\gamma(\Phi^0)^{1/\epsilon}.                       \label{eq:invU}
\end{equation}
Equations \eqref{eq:invwage}, \eqref{eq:income},
\eqref{eq:invLM}--\eqref{eq:invU} cover every primitive and adjusted object
saved by the inversion.

\begin{algorithm}[ht]
\caption{Baseline inversion}
\begin{algorithmic}[1]
\Require $R^{obs},H_M^{obs},H_R^{obs},K,\tau^0$ and parameters
\State Rescale employment so both observed margins sum to $H$
\State Initialize positive $A$ and $B$
\While{either employment margin fails \eqref{eq:invmargins}}
 \State Compute wages from \eqref{eq:invwage}
 \State Compute $\Phi_{ij}$, $\pi_{ij}$, $\widehat H_M$, and $\widehat H_R$
 \State Update $A$ and $B$ using \eqref{eq:abupdate} and damping
 \State Normalize $A$ and scale $B$ consistently with the utility normalization
\EndWhile
\State Compute conditional commuting probabilities and resident income
\State Recover $L_M,L_R,L,\theta,\varphi$ from \eqref{eq:invLM}--\eqref{eq:invvarphi}
\State Recover $a,b$ from \eqref{eq:invfund} and $\bar U$ from \eqref{eq:invU}
\Ensure Baseline primitives, adjusted objects, outcomes, and diagnostics
\end{algorithmic}
\end{algorithm}

\section{Counterfactual equilibrium}
A pure transport experiment changes $\tau^0$ to $\tau^1$ while holding
$\{a,b,\varphi,K\}$ and all parameters fixed. The user may instead combine the
transport change with multiplicative local changes
$\widehat a_i$, $\widehat b_i$, and $\widehat\varphi_i$:
\begin{equation}
 a_i^1=a_i^0\widehat a_i,\qquad
 b_i^1=b_i^0\widehat b_i,\qquad
 \varphi_i^1=\varphi_i^0\widehat\varphi_i.             \label{eq:primitivehats}
\end{equation}
A hat of one leaves that fundamental unchanged; for example, $1.10$ raises it
by ten percent. Conditional on target variables, equations
\eqref{eq:phi}--\eqref{eq:spillovers} explicitly construct all other outcomes.

The optional user input is a polygon shapefile or GeoPackage in arbitrary
spatial units. It need not match the endogenous model grid. Let $G_i$ denote a
target cell, $P_s$ a source policy polygon, and $|\cdot|$ polygon area in a
projected coordinate system. GRID maps any source hat $\widehat x_s$ into
\begin{equation}
 \widehat x_i=1+\sum_s
 \frac{|G_i\cap P_s|}{|G_i|}(\widehat x_s-1).          \label{eq:shockinterpolation}
\end{equation}
Uncovered target area therefore contributes the neutral value one. When source
polygons form a complete non-overlapping partition, equation
\eqref{eq:shockinterpolation} is the standard intersection-area-weighted mean.
Overlapping policy polygons are rejected to prevent double counting. The
standardized model-order table contains \texttt{productivity\_hat},
\texttt{amenity\_hat}, and \texttt{structural\_density\_hat}. Before MATLAB
runs, the reporting layer maps all three interpolated percentage changes,
overlays the original policy-polygon boundaries and transport innovation, and
saves both PDF and PNG versions. This provides a direct visual check of
\eqref{eq:shockinterpolation} and the transport intervention. If no policy
geography is configured, all three hats are one.

The current one-rent specification has no regulatory-wedge primitive; a wedge
shock is therefore rejected rather than silently ignored. Adding one would
require specifying how the wedge enters developer returns or land-use arbitrage
and adding the corresponding inversion and equilibrium conditions.

\subsection{Equations and unknowns}
For $N$ active locations, the closed-city fixed point targets
\begin{equation}
 \underbrace{\{w_i,q_i,Q_i,\theta_i\}_{i=1}^N}_{4N\ {\rm unknowns}}. \label{eq:unknowns}
\end{equation}
There are exactly $4N$ corresponding equilibrium conditions:
\begin{enumerate}
\item $N$ labour first-order conditions,
$w_i=\alpha Y_i/H_{Mi}$;
\item $N$ commercial-rent conditions, using
$q_i=(1-\alpha)Y_i/(\theta_iL_i)$ in commercial locations and
\eqref{eq:mixedrent} in mixed locations;
\item $N$ residential-rent conditions, using
$Q_i=(1-\beta)v_i/[(1-\theta_i)L_i]$ in residential locations and
\eqref{eq:mixedrent} in mixed locations;
\item $N$ land-use conditions: \eqref{eq:mixedshare} in mixed locations and
the appropriate specialization restriction $\theta_i\in\{0,1\}$ elsewhere.
\end{enumerate}
The commuting probabilities, employment margins, $A$, $B$, $Y$, and $v$ are
not omitted equations. They are explicit mappings from the four target vectors
through \eqref{eq:phi}--\eqref{eq:spillovers}; equivalently, one could list
them as additional unknowns together with the same number of defining
equations.

For the open city, total population $H$ is an additional scalar unknown and
\begin{equation}
 U=\bar U                                                     \label{eq:opencondition}
\end{equation}
is the additional scalar equation. The open-city system therefore has
$4N+1$ target unknowns and $4N+1$ target equations.

\subsection{Scenario design: what adjusts and what remains fixed}
The toolkit reports three conceptually different responses to the same
transport intervention. The closed and open cities are market-clearing spatial
equilibria. The fixed-distribution case is a deliberately partial-equilibrium
accounting benchmark. Table~\ref{tab:scenarios} summarizes their identifying
restrictions.

\begin{table}[ht]
\centering\small
\caption{Scenario comparison for a pure transport counterfactual}
\label{tab:scenarios}
\begin{tabular}{P{3.0cm}P{3.4cm}P{3.4cm}P{3.4cm}}
\toprule
 & Closed city & Open city & Fixed distribution\\
\midrule
City population & Fixed at baseline & Adjusts endogenously & Fixed at baseline\\
Expected utility & Adjusts endogenously & Fixed at outside option & Accounting measure adjusts\\
OD assignments & Adjust endogenously & Adjust endogenously & Fixed at baseline\\
Local population and jobs & Adjust endogenously & Adjust endogenously & Fixed at baseline\\
Wages and output & Adjust endogenously & Adjust endogenously & Recomputed at fixed inputs\\
Floor-space rents and land use & Adjust endogenously & Adjust endogenously & Fixed at baseline\\
Transport-mediated spillovers & Recomputed & Recomputed & Recomputed at fixed densities\\
Market clearing & Yes & Yes & No: accounting benchmark\\
\bottomrule
\end{tabular}
\end{table}

The scenarios answer distinct questions. The \emph{closed city} asks how a
fixed set of workers sorts across residences and workplaces once commuting
conditions change. Utility is the natural worker-welfare statistic; GDP and
land rent include the effects of relocation, rent adjustment, and endogenous
spillovers. The \emph{open city} asks how many workers the city attracts when
utility is tied to an outside option. Its near-zero reported utility change is
therefore a model restriction, not an absence of benefits. Benefits are
capitalized into population, output, rents, and land rent. The
\emph{fixed-distribution} case asks what the original commuters would gain if
they retained their baseline OD assignments and local density distributions.
It includes direct commuting gains and travel-time-mediated productivity and
amenity gains, but excludes relocation, resorting, migration, rent
capitalization, and land-use adjustment.

\subsection{Closed city}
The closed city fixes $H=H^0$. The equilibrium mapping is
\begin{align}
 L_i&=\varphi_iK_i^{1-\mu},\\
 \pi_{ij}&=\pi_{ij}(B,Q,w,\tau),&
 (H_R,H_M)&=H\,\mathrm{margins}(\pi),\\
 A_i&=a_i\Upsilon_i(H_M,\tau)^\lambda,&
 B_i&=b_i\Omega_i(H_R,\tau)^\eta,\\
 Y_i&=A_iH_{Mi}^{\alpha}(\theta_iL_i)^{1-\alpha},&
 v_i&=H_{Ri}\sum_j\pi_{j|i}w_j,
\end{align}
followed by the $4N$ target conditions above. Utility from
\eqref{eq:utility} is the aggregate welfare outcome.

\begin{algorithm}[ht]
\caption{Closed-city equilibrium}
\begin{algorithmic}[1]
\Require $a,b,\varphi,K,\tau,H$ and parameters
\State Compute $L=\varphi K^{1-\mu}$ and initialize $w,q,Q,\theta$
\While{any target vector has not converged}
 \State Construct $A$ and $B$ from current employment densities
 \State Construct commuting probabilities and allocate fixed $H$
 \State Recompute $A,B,Y$, conditional probabilities, and $v$
 \State Predict $w,q,Q,\theta$ from the $4N$ equilibrium conditions
 \State Damp each target toward its prediction and record all gaps
\EndWhile
\State Compute $U$ and annual land rent using \eqref{eq:landrent}
\Ensure Equilibrium with fixed $H$ and endogenous $U$
\end{algorithmic}
\end{algorithm}

\subsection{Open city}
The open city fixes $\bar U$ at the inverted baseline level. In addition to the
closed-city mapping, the code uses
\begin{equation}
 H^{*}=H\left(\frac{U}{\bar U}\right)^\epsilon              \label{eq:Hupdate}
\end{equation}
and damps population toward $H^*$. If current city utility exceeds the outside
option, population rises. Convergence requires both all local target equations
and \eqref{eq:opencondition}.

\begin{algorithm}[ht]
\caption{Open-city equilibrium}
\begin{algorithmic}[1]
\Require $a,b,\varphi,K,\tau,\bar U$ and parameters
\State Initialize $H,w,q,Q,\theta$
\While{a local target fails or $U\ne\bar U$}
 \State Execute the closed-city equilibrium mapping for current $H$
 \State Compute $U$ and predict population using \eqref{eq:Hupdate}
 \State Apply a damped population update
\EndWhile
\State Compute annual land rent using \eqref{eq:landrent}
\Ensure Equilibrium with fixed $\bar U$ and endogenous $H$
\end{algorithmic}
\end{algorithm}

\subsection{Fixed-distribution accounting benchmark}
\label{sec:fixeddistribution}
The toolkit also reports a deliberately constrained benchmark in which the
baseline residence--workplace assignment is held fixed. This exercise is
designed to isolate gains available without spatial reallocation. It is
\emph{not} a third market-clearing equilibrium closure.

Let a superscript $F$ denote the fixed-distribution counterfactual. The
benchmark holds fixed
\begin{equation}
 H^F=H^0,\qquad H_{Mi}^F=H_{Mi}^0,\qquad H_{Ri}^F=H_{Ri}^0,
 \qquad \pi_{ij}^F=\pi_{ij}^0.                         \label{eq:fixedmargins}
\end{equation}
Thus neither the number of modeled workers, their workplace and residence
distributions, nor their bilateral origin--destination assignments change.
Residential and commercial floor-space rents and the commercial land-use share
are also held at baseline values:
\begin{equation}
 Q_i^F=Q_i^0,\qquad q_i^F=q_i^0,\qquad \theta_i^F=\theta_i^0. \label{eq:fixedprices}
\end{equation}
For a pure transport experiment, $\varphi$ and $K$ are unchanged, so total,
commercial, and residential floor space are fixed as well.

Fixed local quantities do \emph{not} fix endogenous spillovers. Both effective
density kernels depend on bilateral travel times. The benchmark therefore
recomputes
\begin{align}
 \Upsilon_i^F&=\sum_s e^{-\delta\tau_{is}^1}
                    \frac{H_{Ms}^0}{K_s},
 &A_i^F&=a_i^1(\Upsilon_i^F)^\lambda,\\
 \Omega_i^F&=\sum_r e^{-\rho\tau_{ir}^1}
                    \frac{H_{Rr}^0}{K_r},
 &B_i^F&=b_i^1(\Omega_i^F)^\eta.                      \label{eq:fixedspillovers}
\end{align}
Shorter travel times can consequently raise effective workplace and
residential density even though the underlying employment and population
vectors do not move.

With workplace labour and commercial floor-space inputs fixed, counterfactual
output and the competitive wage are
\begin{equation}
 Y_i^F=A_i^F(H_{Mi}^0)^\alpha(\theta_i^0L_i^0)^{1-\alpha},
 \qquad w_i^F=\frac{\alpha Y_i^F}{H_{Mi}^0}.           \label{eq:fixedoutput}
\end{equation}
Wages are therefore \emph{not} fixed: productivity-spillover gains pass into
output and wages. The benchmark recomputes resident labour income using the
fixed OD probabilities and $w^F$, but it does not use that income to clear the
residential floor-space market because $Q^F=Q^0$ is maintained.

For a worker who remains on baseline route $(i,j)$, baseline floor-space rents
cancel from the proportional indirect-utility change. The log change is
\begin{equation}
 \Delta\log u_{ij}^F=
 \log\frac{B_i^F}{B_i^0}
 +\log\frac{w_j^F}{w_j^0}
 -\kappa(\tau_{ij}^1-\tau_{ij}^0).                    \label{eq:fixedrouteutility}
\end{equation}
The aggregate fixed-distribution log-welfare change is the baseline-flow
weighted average
\begin{equation}
 g^F=\sum_i\sum_j\pi_{ij}^0\Delta\log u_{ij}^F
     =g_B^F+g_w^F+g_\tau^F,                            \label{eq:fixedwelfare}
\end{equation}
where
\begin{align}
 g_B^F&=\sum_{ij}\pi_{ij}^0\log(B_i^F/B_i^0),\\
 g_w^F&=\sum_{ij}\pi_{ij}^0\log(w_j^F/w_j^0),\\
 g_\tau^F&=-\kappa\sum_{ij}\pi_{ij}^0
                 (\tau_{ij}^1-\tau_{ij}^0).           \label{eq:fixedcomponents}
\end{align}
The table reports total worker welfare as $100[\exp(g^F)-1]$ and each component
as $100[\exp(g_k^F)-1]$. Because the components combine additively in logs but
multiplicatively in percentage levels, the displayed component percentages
need not sum exactly to the displayed total. Fixed-distribution GDP is
\begin{equation}
 100\left(\frac{\sum_iY_i^F}{\sum_iY_i^0}-1\right).    \label{eq:fixedgdp}
\end{equation}

For a pure transport experiment, annual land rent is unchanged by construction:
$q$, $Q$, $\theta$, and $L$ are all fixed in \eqref{eq:landrent}. Productivity
gains accrue to wages in \eqref{eq:fixedoutput}; amenity and commuting gains
accrue to workers in \eqref{eq:fixedrouteutility}. The zero land-rent change is
therefore an assumption of this partial-equilibrium benchmark, not a prediction
that landowners cannot benefit in the closed- or open-city equilibria. Allowing
rents to adjust while imposing \eqref{eq:fixedmargins} would generally violate
one or more housing-market, firm, and mixed-use land-arbitrage conditions unless
additional location-specific wedges were introduced.

\subsection{Zero-spillover versions of all three scenarios}
When \texttt{project.run\_no\_spillover\_comparison} is enabled, the toolkit
also runs closed-city, open-city, and fixed-distribution versions with
\begin{equation}
 \lambda=0,\qquad \eta=0.                              \label{eq:nospillovers}
\end{equation}
This switch removes the endogenous productivity response to employment
effective density and the endogenous amenity response to residential effective
density. The decay parameters $\delta$ and $\rho$ remain in the parameter
structure but are economically inactive because the corresponding
elasticities are zero. Fundamentals, transport times, preferences, production,
floor-space supply, commuting choice, and the applicable closure condition
remain in the model.

The zero-spillover run is not obtained by merely setting two parameters to zero
after loading the main inversion. It re-inverts the observed baseline under
\eqref{eq:nospillovers} and then solves the baseline and counterfactual under
that same specification. This makes each reported percentage change an
internally consistent comparison and ensures that differences between the two
specifications reflect spillovers rather than mismatched baseline primitives.
The main results are unchanged and remain in \path{outputs/simulation}. The
sensitivity outputs are stored separately in
\path{outputs/no_spillovers/simulation}; the combined comparison is written to
\path{outputs/simulation/aggregate_changes_spillover_comparison.csv}.

Interpretation differs by scenario:
\begin{itemize}
\item In the \textbf{closed city without spillovers}, population remains fixed,
but workers may relocate, wages, rents, and land use may adjust, and utility
changes. The difference from the main closed-city result measures the
contribution of endogenous productivity and amenity feedbacks, including their
interaction with relocation.
\item In the \textbf{open city without spillovers}, utility remains fixed and
population adjusts. Migration, commuting resorting, wages, rents, land use,
GDP, and land rent remain endogenous. The difference from the main open-city
result shows how spillovers affect the city's population and capitalization
response.
\item In the \textbf{fixed distribution without spillovers}, OD assignments,
population, employment, rents, and land use remain fixed. With no direct
productivity or amenity shock, productivity, amenities, wages, GDP, and land
rent therefore remain unchanged. Worker welfare changes only through the
commuting-cost term $g_\tau^F$. This is the cleanest model-based measure of the
benefit of shorter travel times for baseline commuters.
\end{itemize}

The difference between the with- and without-spillover results is a sensitivity
comparison, not in general an additive ``spillover component.'' In the two
equilibrium closures, spillovers change sorting, prices, land use, and (in the
open city) population, so interactions are part of the difference. Only the
fixed-distribution decomposition in \eqref{eq:fixedwelfare} cleanly separates
commuting, productivity/wage, and amenity terms in log welfare.

\section{Parameter choices}
\begin{table}[ht]
\centering\small
\caption{Current demonstration defaults in \texttt{project\_config.yaml}}
\begin{tabular}{llll}
\toprule Parameter & Configuration key & Default & Interpretation\\
\midrule
$\alpha$ & \texttt{production\_share\_labor} & 0.80 & Labour share in production\\
$\beta$ & \texttt{expenditure\_share\_consumption} & 0.75 & Numeraire expenditure share\\
$\epsilon$ & \texttt{commuting\_elasticity} & 6.6941 & Frechet shape/commuting elasticity\\
$\kappa$ & \texttt{commuting\_time\_coefficient} & 0.014744 & Cost per travel-time unit\\
$\lambda$ & \texttt{productivity\_spillover} & 0.0710 & Productivity feedback\\
$\delta$ & \texttt{productivity\_spatial\_decay} & 0.3617 & Productivity decay per time unit\\
$\eta$ & \texttt{amenity\_spillover} & 0.1553 & Amenity feedback\\
$\rho$ & \texttt{amenity\_spatial\_decay} & 0.7595 & Amenity decay per time unit\\
$\mu$ & \texttt{construction\_capital\_share} & 0.75 & Non-land development share\\
$1-\mu$ & \texttt{construction\_land\_share} & 0.25 & Land share\\
\bottomrule
\end{tabular}
\end{table}

These defaults reproduce the ARSW-style demonstration parameterization; they
are not universal constants. For a new city, the user should document the
source and empirical setting for every parameter. In particular:
\begin{itemize}
\item $\kappa$, $\delta$, and $\rho$ must be compatible with the unit of the
travel-time matrices; changing minutes to hours without rescaling them changes
the model.
\item $\epsilon$ governs both sorting responses and the open-city population
update, so welfare and migration results can be sensitive to it.
\item $\lambda$ and $\eta$ govern amplification through endogenous productivity
and amenity. Zero is the transparent no-spillover sensitivity case.
\item $\alpha$, $\beta$, and $1-\mu$ determine the incidence of revenue across
labour, floor-space uses, and land. In particular, annual land rent in
\eqref{eq:landrent} uses $1-\mu=0.25$.
\item Numerical tolerances and damping weights are algorithmic controls, not
economic parameters; changing them should affect convergence, not the
converged equilibrium.
\end{itemize}

\section{Counterfactual outputs and validation}
For each requested closure, the toolkit solves both $\tau^0$ and $\tau^1$,
computes block-level percentage changes, and reports aggregate utility,
population, GDP, and annual land rent. The main aggregate table uses
\begin{align}
 \%\Delta GDP &=100\left(\frac{\sum_iY_i^1}{\sum_iY_i^0}-1\right),\\
 \%\Delta\Rland&=100\left(
 \frac{\sum_i(1-\mu)[q_i^1\theta_i^1L_i+
 Q_i^1(1-\theta_i^1)L_i]}
 {\sum_i(1-\mu)[q_i^0\theta_i^0L_i+
 Q_i^0(1-\theta_i^0)L_i]}-1\right).                    \label{eq:reportland}
\end{align}
For a pure transport counterfactual $\varphi$ and $K$ are fixed, hence $L$ is
fixed, but $q$, $Q$, and $\theta$ need not be. Consequently, percentage changes
in annual land rent are not generally equal to changes in either floor-space
rent.

\subsection{Three travel-time measures}
Let $\pi^0$ and $\pi^1$ denote baseline and counterfactual unconditional OD
probabilities, $\tau^0$ and $\tau^1$ the corresponding travel-time matrices,
and $H^0,H^1$ modeled commuter populations. Define
\begin{equation}
 \bar T_{00}=\sum_{ij}\pi_{ij}^0\tau_{ij}^0,\qquad
 \bar T_{01}=\sum_{ij}\pi_{ij}^0\tau_{ij}^1,\qquad
 \bar T_{11}=\sum_{ij}\pi_{ij}^1\tau_{ij}^1.          \label{eq:traveltimemeans}
\end{equation}
The aggregate output table distinguishes:
\begin{description}
\item[Mean commute time, baseline flows]
\path{MeanCommuteTimePct_BaselineFlows} is
$100(\bar T_{01}/\bar T_{00}-1)$. It changes the network while retaining the
baseline OD weights. It is the immediate or before-relocation network effect
for the original commuting pattern.
\item[Mean commute time, counterfactual flows]
\path{MeanCommuteTimePct_CounterfactualFlows} is
$100(\bar T_{11}/\bar T_{00}-1)$. It uses post-adjustment OD weights and hence
includes the effect of changed residence--workplace patterns on the average
commute.
\item[Aggregate commuter-minutes]
\path{AggregateCommuterMinutesPct} is
$100[H^1\bar T_{11}/(H^0\bar T_{00})-1]$. It measures the total one-way minutes
across all modeled commuters. It equals the counterfactual-flow mean measure
when population is fixed, but differs in the open city because $H^1$ changes.
\end{description}
All three measures use one-way commuting time. They are physical time
statistics, not monetized benefits. The fixed-distribution benchmark has
$\pi^1=\pi^0$ and $H^1=H^0$, so all three percentage changes coincide.

Before interpretation, users should verify:
\begin{enumerate}
\item location IDs and matrix ordering are identical;
\item employment and population aggregate to the same number;
\item inversion reproduces both observed employment margins;
\item all fixed-point gaps satisfy the configured tolerance;
\item baseline equilibrium reproduction is satisfactory;
\item land-rent output is constructed from \eqref{eq:landrent}, not from an
input field labelled land value.
\end{enumerate}

\section{Code correspondence}
\begin{longtable}{P{4.3cm}P{9.2cm}}
\toprule File & Economic role\\ \midrule
\texttt{invert\_baseline.m} & Commented baseline-inversion entry script.\\
\texttt{run\_counterfactual.m} & Solves the requested equilibrium closures and the fixed-distribution benchmark.\\
\texttt{run\_no\_spillovers.m} & Re-inverts and reruns all three scenarios with $\lambda=\eta=0$, saving results separately.\\
\texttt{qt\_invert\_baseline.m} & Passes the single observed rent into both use-specific inversion conditions.\\
\texttt{cmodexog.m} & Inverts $A,B,w$ and commuting probabilities jointly.\\
\texttt{cdensityE.m} & Inverts $L_M,L_R,L,\theta,\varphi$.\\
\texttt{cprod.m}, \texttt{cres.m} & Recover $a$ and $b$ from total productivity and amenity.\\
\texttt{ubar.m} & Computes baseline expected utility $\bar U$.\\
\texttt{smodendog.m} & Solves the closed-city $4N$-equation target system.\\
\texttt{ussmodendog.m} & Adds population and utility clearing for the open city.\\
\path{qt_solve_fixed_distribution.m} & Evaluates fixed OD assignments with transport-mediated spillovers and decomposes worker welfare.\\
\path{qt_write_closure_results.m} & Reports local outcomes, annual land rent, and the three travel-time measures.\\
\path{qt_display_results.m} & Displays compact MATLAB headings while preserving descriptive CSV variable names.\\
\bottomrule
\end{longtable}

\section{Interpretation and limitations}
Results are comparative statics conditional on the model, parameters, baseline
data, and travel times. Wages are always recovered inside the model inversion;
they are never generated from an auxiliary employment elasticity. Synthetic
rents are useful for testing the pipeline but not for policy valuation. Observed rents require compatible
units, concepts, and dates. Annual land rent is a model-implied flow; converting
it to a land asset value requires an explicit discounting and expectations
model. Open-city utility is fixed by construction, so population rather than
utility is its aggregate adjustment margin.

\section*{Reference}
Ahlfeldt, G. M., Redding, S. J., Sturm, D. M., and Wolf, N. (2015).
``The Economics of Density: Evidence from the Berlin Wall.''
\emph{Econometrica}, 83(6), 2127--2189.
\href{https://doi.org/10.3982/ECTA10876}{doi:10.3982/ECTA10876}.
\end{document}
